\section{The chamber torsor inside the nine-dimensional matrix carrier}
\label{rnmt:sec:matrix-torsor}

The preceding section separates the two three-element basis orbits of
\(\Lambda_{\rm tor}\) from the regular six-element set of Kuhn chambers.
There is nevertheless a natural place where the regular six-set occurs.  It
is not the lattice \(\Lambda_{\rm tor}\) itself, but the off-diagonal part of
the rank-nine bilinear carrier
\[
 \mathcal M_{\mathbb Z}
 :=\operatorname{Hom}(\overline L_3,L_3)
 \cong\operatorname{Mat}_{3\times3}(\mathbb Z).
\]
This is the integral real form underlying the nine-real-dimensional
Hermitian carrier.  The distinction between the full carrier and its
positive-semidefinite cone will be essential.

Retain the marked bases \(e_1,e_2,e_3\) of \(L_3\) and
\(\bar e_1,\bar e_2,\bar e_3\) of \(\overline L_3\).  Let \(E_{ij}\) denote
the homomorphism which sends \(\bar e_j\) to \(e_i\) and annihilates the
other two basis vectors.  The diagonal \(S_3\)-action on the two rank-three
lattices induces
\begin{equation}
 \sigma\cdot X=R_\sigma X R_\sigma^{-1},
 \qquad \sigma\cdot E_{ij}=E_{\sigma(i),\sigma(j)}.
 \label{rnmt:eq:conjugation-action}
\end{equation}
Put
\begin{equation}
 \Delta=\bigoplus_{i=1}^3\mathbb ZE_{ii},
 \qquad
 \mathcal O=\bigoplus_{i\ne j}\mathbb ZE_{ij}.
 \label{rnmt:eq:diagonal-offdiagonal}
\end{equation}

\begin{theorem}[The regular chamber torsor in rank nine]
\label{rnmt:thm:regular-torsor}
There is an integral \(S_3\)-module decomposition
\begin{equation}
 \mathcal M_{\mathbb Z}=\Delta\oplus\mathcal O,
 \qquad
 \Delta\cong\mathbb Z[S_3/S_2],
 \qquad
 \mathcal O\cong\mathbb Z[S_3].
 \label{rnmt:eq:module-decomposition}
\end{equation}
If a chamber is indexed in one-line notation by
\(\pi=(\pi_1\pi_2\pi_3)\in S_3\), then
\begin{equation}
 \vartheta(c_\pi)=E_{\pi_1,\pi_3}
 \label{rnmt:eq:chamber-matrix-bijection}
\end{equation}
is an \(S_3\)-equivariant bijection from the six Kuhn chambers at a marked
vertex to the six-element basis of \(\mathcal O\).  Equivalently, these are
the six ordered mixed pairs
\begin{equation}
 \bigl((e_i,0),(0,\bar e_j)\bigr),\qquad i\ne j.
 \label{rnmt:eq:offdiagonal-pairs}
\end{equation}
They are exactly the mixed basis pairs which are orthogonal for the integral
alternating form \(E_3\), since
\begin{equation}
 E_3\bigl((e_i,0),(0,\bar e_j)\bigr)=\delta_{ij}.
 \label{rnmt:eq:mixed-pairing}
\end{equation}
The three excluded diagonal pairs \(i=j\) comprise the basis of \(\Delta\)
and are the symplectically dual mixed pairs.
\end{theorem}

\begin{proof}
Equation \eqref{rnmt:eq:diagonal-offdiagonal} partitions the nine matrix
units and is preserved by \eqref{rnmt:eq:conjugation-action}, proving the
direct-sum assertion.  The action is transitive on the three diagonal units,
and the stabilizer of \(E_{11}\) is the subgroup permuting \(2\) and \(3\);
hence \(\Delta\cong\mathbb Z[S_3/S_2]\).

The action on ordered pairs \((i,j)\) with \(i\ne j\) is free: a permutation
fixing two distinct letters fixes the remaining letter and is therefore the
identity.  It is transitive because any ordered pair of distinct letters can
be carried to any other.  Thus \(\mathcal O\) is the regular permutation
module.  For each ordered pair \((i,j)\), there is a unique remaining letter
\(k\), and hence a unique one-line permutation \((i,k,j)\).  This proves
that \eqref{rnmt:eq:chamber-matrix-bijection} is a bijection.  Moreover,
\[
 \vartheta(c_{\sigma\pi})
 =E_{\sigma(\pi_1),\sigma(\pi_3)}
 =\sigma\cdot E_{\pi_1,\pi_3},
\]
so it is equivariant.  Finally, \eqref{rnmt:eq:mixed-pairing} is immediate
from the defining formula
\(E_3((a,b),(a',b'))=a^{\mathsf T}b'-b^{\mathsf T}a'\).
\end{proof}

\subsection{The signed quotient and the exact chamber incidence}
\label{rnmt:subsec:signed-quotient}

The repeated indices on the angular and deck bases determine a marking
\(\iota:\overline L_3\to L_3\), \(\iota(\bar e_i)=e_i\).  Define
\begin{equation}
 s:A_F\oplus B_F\longrightarrow A_F,\qquad
 s(x,y)=x-B_0y,
 \qquad
 D=s\circ(P\oplus Q).
 \label{rnmt:eq:D-factorization}
\end{equation}
Since \(B_0Q=P\), this means, in the retained bases,
\begin{equation}
 D(a,b)=Pa-B_0Qb=P(a-\iota b),\qquad [D]=[P\mid-P].
 \label{rnmt:eq:D-formula}
\end{equation}
The formula is typed: \(Q\) is applied to the deck lattice and \(B_0\)
then carries its value from \(B_F\) to \(A_F\).

\begin{theorem}[Unimodular pair--Kuhn--Weyl incidence]
\label{rnmt:thm:pair-kuhn-weyl}
The map \(D\) is surjective, with saturated kernel
\begin{equation}
 \ker D
 =\mathbb Z\langle
 (e_1,\bar e_1),(e_2,\bar e_2),(e_3,\bar e_3),(k,0)\rangle,
 \qquad k=e_1+e_2+e_3.
 \label{rnmt:eq:D-kernel}
\end{equation}
For \(i\ne j\), put
\[
 U_{ij}=\mathbb Z(e_i,0)\oplus\mathbb Z(0,\bar e_j).
\]
Then \(U_{ij}\) is primitive and isotropic for \(E_3\), and
\begin{equation}
 D|_{U_{ij}}:U_{ij}\xrightarrow{\ \sim\ }A_F,\qquad
 (e_i,0)\longmapsto n_i,\quad(0,\bar e_j)\longmapsto-n_j
 \label{rnmt:eq:unimodular-plane}
\end{equation}
is an integral isomorphism.

Let \(k_{ij}\) be the remaining index and
\(\pi_{ij}=(i,k_{ij},j)\).  Under the Cartan marking
\(\jmath(n_i)=\nu_i=E_{ii}-I_3/3\), the positive quadrant in \(U_{ij,\mathbb
R}\) maps isomorphically to the Weyl chamber
\begin{equation}
 \mathcal C_{ij}
 =\mathbb R_{\geq0}\nu_i+\mathbb R_{\geq0}(-\nu_j)
 =\{\operatorname{diag}(\lambda):
 \lambda_i\geq\lambda_{k_{ij}}\geq\lambda_j,\
 \textstyle\sum_r\lambda_r=0\}.
 \label{rnmt:eq:chamber-inequalities}
\end{equation}
The assignments
\begin{equation}
 ((e_i,0),(0,\bar e_j))
 \longmapsto \pi_{ij}
 \longmapsto[\sigma_{0,\pi_{ij}}]
 \longmapsto\mathcal C_{ij}
 \label{rnmt:eq:torsor-chain}
\end{equation}
are equivariant bijections of \(S_3\)-torsors, and
\begin{equation}
 P(\sigma_{0,\pi_{ij}})
 =\operatorname{conv}\{0,n_i,-n_j\}.
 \label{rnmt:eq:projected-kuhn}
\end{equation}
The vector
\begin{equation}
 D(e_i,\bar e_j)=n_i-n_j
 \label{rnmt:eq:interior-root}
\end{equation}
is the interior \(A_2\)-root bisector of this chamber.  With the source and
target orientations fixed by the displayed bases,
\begin{equation}
 \det[e_i\mid e_i+e_{k_{ij}}\mid k]
 =\operatorname{sgn}(\pi_{ij})
 =\det[n_i\mid -n_j].
 \label{rnmt:eq:orientation-sign}
\end{equation}
\end{theorem}

\begin{proof}
Since \(P\) is surjective, so is \(D\).  Equation
\eqref{rnmt:eq:D-formula} gives \(D(a,b)=0\) precisely when
\(a-\iota b=tk\) for a unique \(t\in\mathbb Z\).  Hence
\[
 (a,b)=\sum_{r=1}^3 b_r(e_r,\bar e_r)+t(k,0),
\]
and uniqueness gives \eqref{rnmt:eq:D-kernel}.  Surjectivity implies that
the kernel is saturated.

The plane \(U_{ij}\) is spanned by two members of a basis of
\(\Lambda_{\rm tor}\), hence is primitive.  Equation
\eqref{rnmt:eq:mixed-pairing} makes it isotropic exactly when \(i\ne j\).
In its displayed basis the restriction of \(D\) has columns
\([n_i\mid-n_j]\).  Directly from
\[
 n_1=\binom10,\qquad n_2=\binom{-1}{1},\qquad
 n_3=\binom0{-1},
\]
one obtains
\(\det[n_i\mid-n_j]=\pm1\) for \(i\ne j\), proving
\eqref{rnmt:eq:unimodular-plane}.

For \(x,y\geq0\), the diagonal matrix
\(H=x\nu_i-y\nu_j\) satisfies
\[
 \lambda_i-\lambda_{k_{ij}}=x,\qquad
 \lambda_{k_{ij}}-\lambda_j=y.
\]
Conversely these two differences, together with trace zero, recover \(H\),
which proves \eqref{rnmt:eq:chamber-inequalities} and the cone
isomorphism.  The four vertices of
\(\sigma_{0,(i,k_{ij},j)}\) have \(P\)-images
\[
 0,\quad n_i,\quad n_i+n_{k_{ij}}=-n_j,\quad
 n_1+n_2+n_3=0,
\]
proving \eqref{rnmt:eq:projected-kuhn}.  Ordered distinct pairs, one-line
permutations, Kuhn-cell translation classes, and Weyl chambers are all free
transitive \(S_3\)-sets, and the displayed assignments commute with the
coordinate action.  At \(x=y=1\),
\(\jmath(n_i-n_j)=E_{ii}-E_{jj}\), whose ordered diagonal entries in
positions \((i,k_{ij},j)\) are \((1,0,-1)\); this is the interior bisector.

Finally, relative to the ordered basis
\((e_i,e_{k_{ij}},e_j)\), the columns
\((e_i,e_i+e_{k_{ij}},k)\) form an upper unitriangular matrix.  Their
determinant in the fixed source basis is therefore
\(\operatorname{sgn}(\pi_{ij})\).  For \((i,k_{ij},j)=(1,2,3)\),
\([n_i\mid-n_j]=I_2\); equivariance and
\(\det\rho_A(\sigma)=\operatorname{sgn}(\sigma)\) prove the second equality
in \eqref{rnmt:eq:orientation-sign}.
\end{proof}

\begin{proposition}[Exact marking and monodromy boundary]
\label{rnmt:prop:marking-boundary}
There is no endomorphism \(T:A_F\to A_F\) satisfying
\begin{equation}
 DM_3=TD.
 \label{rnmt:eq:no-monodromy-intertwiner}
\end{equation}
Moreover, for every \(c\in\mathbb Z\), the shear
\begin{equation}
 S_c(a,b)=(a+c\iota b,b)
 \label{rnmt:eq:marking-shear}
\end{equation}
is an integral automorphism which fixes the angular lattice and its Kuhn
fan, preserves \(c_3\) and \(E_3\), and commutes with \(M_3\) and the
diagonal \(S_3\)-action, but
\begin{equation}
 DS_c(a,b)=D(a,b)+cP(\iota b).
 \label{rnmt:eq:D-shear}
\end{equation}
Thus the pair--chamber correspondence is canonical relative to the retained
angular/deck bases, their common indexing, \(P,Q,B_0\), and the sign in
\eqref{rnmt:eq:D-factorization}; it is not an invariant of the unmarked data
\((M_3,E_3,c_3)\) together with the angular fan.
\end{proposition}

\begin{proof}
The vector \((e_i,\bar e_i)\) lies in \(\ker D\), whereas
\[
 M_3(e_i,\bar e_i)=(2e_i,\bar e_i),\qquad
 D(2e_i,\bar e_i)=n_i\ne0.
\]
Thus \(\ker D\) is not \(M_3\)-invariant, which is necessary for
\eqref{rnmt:eq:no-monodromy-intertwiner}.

The inverse of \(S_c\) is \(S_{-c}\).  Direct substitution gives
\[
 S_cM_3(a,b)=(a+(c+1)\iota b,b)=M_3S_c(a,b),
\]
and proves preservation of \(c_3\) and commutation with the diagonal
coordinate action.  Symmetry of the dot product gives
\[
 E_3(S_c(a,b),S_c(a',b'))
 =(a+c\iota b)^{\mathsf T}b'
   -b^{\mathsf T}(a'+c\iota b')
 =E_3((a,b),(a',b')).
\]
Equation \eqref{rnmt:eq:D-shear} follows from
\eqref{rnmt:eq:D-formula}.  For \(c\ne0\), the deck complement is moved, so
the retained marked pair set and \(D\) are changed.
\end{proof}

The parity and opposition operations also retain exact information.  The
even and odd chamber triples are
\begin{align}
 \vartheta\{c_{123},c_{231},c_{312}\}
   &=\{E_{13},E_{21},E_{32}\},
 \label{rnmt:eq:even-cycle}\\
 \vartheta\{c_{132},c_{213},c_{321}\}
   &=\{E_{12},E_{23},E_{31}\}.
 \label{rnmt:eq:odd-cycle}
\end{align}
Thus they are the two orientations of the directed three-cycle.  Reversal of
one-line notation, \((i,k,j)\mapsto(j,k,i)\), is carried by \(\vartheta\) to
matrix transposition \(E_{ij}\mapsto E_{ji}\).  This is the exact chamber
opposition involution; it exchanges the two parity triples.

\subsection{The real Hermitian representation and its integral defect}

Let \(\operatorname{Herm}_3(\mathbb C)\) carry the \(S_3\)-action by
conjugation with the real permutation matrices.  Define
\begin{equation}
 \Xi:\operatorname{Mat}_{3\times3}(\mathbb R)
       \longrightarrow\operatorname{Herm}_3(\mathbb C),
 \qquad
 \Xi(X)=\frac{X+X^{\mathsf T}}2
       +\frac{i(X-X^{\mathsf T})}2.
 \label{rnmt:eq:Xi}
\end{equation}

\begin{proposition}[Exact representation bridge and the positivity boundary]
\label{rnmt:prop:representation-bridge}
The map \(\Xi\) is an \(S_3\)-equivariant real-linear isomorphism.  Its
inverse sends a Hermitian matrix \(H=S+iA\), with \(S\) real symmetric and
\(A\) real skew-symmetric, to \(S+A\).  Consequently
\begin{equation}
 \operatorname{Herm}_3(\mathbb C)
 \cong_{\mathbb R[S_3]}
 2\mathbin{\cdot}\mathbf 1\oplus
 3\mathbin{\cdot}V_{\rm std}\oplus\mathbf{sgn}.
 \label{rnmt:eq:hermitian-character}
\end{equation}

For each unordered pair \(\{i,j\}\), the symmetric and skew generators
\(E_{ij}+E_{ji}\) and \(E_{ij}-E_{ji}\) span an index-two sublattice of
\(\mathbb ZE_{ij}\oplus\mathbb ZE_{ji}\).  Hence the symmetric--skew
off-diagonal splitting has total integral index
\begin{equation}
 2^3=8.
 \label{rnmt:eq:index-eight}
\end{equation}
In particular, \(\Xi\) is a real representation isomorphism, not an
identification of the two preferred integral lattices.

Moreover, for \(i\ne j\), \(\Xi(E_{ij})\) has eigenvalues
\(2^{-1/2},-2^{-1/2},0\).  Thus the regular chamber torsor is located in the
nine-dimensional Hermitian representation, but its matrix-unit image is not
contained in the positive-semidefinite Gram cone or its boundary.
\end{proposition}

\begin{proof}
Every real matrix has the unique decomposition
\(X=S+A\), where
\(S=(X+X^{\mathsf T})/2\) is symmetric and
\(A=(X-X^{\mathsf T})/2\) is skew-symmetric.  Equation
\eqref{rnmt:eq:Xi} sends this pair to \(S+iA\), which is Hermitian, and the
displayed inverse proves bijectivity.  Permutation conjugation commutes with
transpose and with these two projections, proving equivariance.

The real permutation representation on three letters is
\(\mathbf1\oplus V_{\rm std}\).  Therefore
\[
 \operatorname{Mat}_{3\times3}(\mathbb R)
 \cong(\mathbf1\oplus V_{\rm std})^{\otimes2}
 \cong2\mathbin{\cdot}\mathbf 1\oplus
       3\mathbin{\cdot}V_{\rm std}\oplus\mathbf{sgn},
\]
where
\(V_{\rm std}\otimes V_{\rm std}
 \cong\mathbf1\oplus V_{\rm std}\oplus\mathbf{sgn}\).
For one unordered pair, the change-of-basis matrix from
\((E_{ij},E_{ji})\) to
\((E_{ij}+E_{ji},E_{ij}-E_{ji})\) has determinant \(-2\).  Multiplication
over the three unordered pairs proves \eqref{rnmt:eq:index-eight}.

Finally, on the \(i,j\) coordinate plane, \(\Xi(E_{ij})\) has block
\[
 \begin{pmatrix}0&(1+i)/2\\(1-i)/2&0\end{pmatrix},
\]
whose eigenvalues are \(\pm |1+i|/2=\pm2^{-1/2}\); the remaining coordinate
has eigenvalue zero.  Hence it is indefinite.
\end{proof}

\subsection{The positive six-set is a different equivariant realization}

Although the matrix units themselves are indefinite, the spectral-floor map
of the preceding section gives a second realization of the same chamber
torsor in the rank-two positive-semidefinite boundary.  If
\(\pi=(i,k,j)\), then
\begin{equation}
 G_\pi=2E_{ii}+E_{kk}
 \label{rnmt:eq:positive-six-set}
\end{equation}
has diagonal entries equal to a permutation of \((2,1,0)\).  The assignment
\begin{equation}
 E_{ij}\longmapsto 2E_{ii}+E_{kk},
 \qquad \{i,j,k\}=\{1,2,3\},
 \label{rnmt:eq:offdiagonal-positive-map}
\end{equation}
is therefore an equivariant bijection of the two marked six-element sets.
It extends uniquely by linearity to a homomorphism
\(\mathcal O\to\Delta\), but that extension is neither the map \(\Xi\) nor
the nonlinear spectral-floor map on all of
\(\operatorname{Herm}_3^0(\mathbb C)\).  Thus
\eqref{rnmt:eq:offdiagonal-positive-map} records exact chamber incidence; it
does not manufacture a positive monodromy-compatible morphism from the
toroidal family.

\begin{auditrecord}[What the native \(\mathbb Z^6\) result becomes]
The six Kuhn chambers are not the six generators of
\(H_1(T^6;\mathbb Z)\).  They are canonically, relative to the retained
marking, the six symplectically orthogonal ordered angular--deck pairs and
hence the six off-diagonal basis elements in the associated rank-nine
bilinear carrier.  This places the regular \(S_3\)-torsor exactly inside the
nine-dimensional program while preserving the different roles of angular
cycles, deck cycles, monodromy, and the positive-semidefinite Gram boundary.
The construction is a representation and incidence theorem; positivity and
global toroidal descent remain separate questions and are not inferred from
the equality of dimensions or orbit cardinalities.
\end{auditrecord}
